% P:p2-s-contents
This supplement gives the amine-specific parameter values at full precision (Section~\ref{sec:p2-s-parameters}), the source conventions of the data (Section~\ref{sec:p2-s-observations}), the deviations inside and outside the measured temperature ranges of the equilibrium constants (Section~\ref{sec:p2-s-ranges}), the published comparisons on the same data (Section~\ref{sec:p2-s-published}), the stability penalty and retreat procedure (Section~\ref{sec:p2-s-retreat}) and the phase-search conditions and reproducibility records (Section~\ref{sec:p2-s-numerical}).

\section{Model inputs and parameter values}\label{sec:p2-s-parameters}
% P:p2-s-parameters
Table~\ref{tab:p2-dea-all-current} lists the value, unit and fitted or fixed role of each amine-specific DEA parameter; the archived parameter file adds the bounds, conditional standard error and source of each.
For the carbamate hydrolysis (R4) and DEA dissociation (R5) correlations, the coefficients give $\ln K(T)=a+b/T+dT$ on the common-molality reference basis, with $T$ in K; the listed coefficients $c$, $e$ and $f$ are zero.
The fitted-parameter table of the main text uses increments relative to the final R4 correlation; the final coefficients below already include the fitted enthalpy change.
Fixed ionic parameters are inherited model assumptions, not independent measurements of ionic size.
The fitted MDEA and AMP parameters and the parameters of their species are in Tables~\ref{main-tab:p2-mdea-species-current}, \ref{main-tab:p2-mdea-fitted-current}, \ref{main-tab:p2-amp-species-current} and \ref{main-tab:p2-amp-fitted-current} of the main text; the archived parameter files give every fixed and fitted value at full precision.
Packing diameter is the hard-sphere diameter used in the repulsive contribution; dielectric volume and dielectric temperature are the coefficients $d_V$ and $d_T$ in the relative-permittivity expression transferred from Part I \cite[Sect.~2.1]{PolleyPartI}.

\begingroup\small\setlength{\tabcolsep}{2pt}\begin{longtable}{@{}>{\raggedright\arraybackslash}p{0.46\linewidth}>{\raggedright\arraybackslash}p{0.21\linewidth}>{\raggedright\arraybackslash}p{0.19\linewidth}>{\raggedright\arraybackslash}p{0.09\linewidth}@{}}
\caption{Complete DEA parameter values. Fitted-coordinate errors and bounds are in the main text; reaction coefficients are completed coefficients, not additional independent fitted quantities. The electronic parameter table retains each value's source and basis.}\label{tab:p2-dea-all-current}\\
\toprule Quantity & Value & Unit & Role \\ \midrule
\endfirsthead\toprule Quantity & Value & Unit & Role \\ \midrule
\endhead\bottomrule\endfoot
$k_{\mathrm{DEAH^+,w}}$ & $-0.18825647$ & 1 & fitted \\
$k_{\mathrm{DEACOO^-,w}}$ & $-0.038548699$ & 1 & fitted \\
hydrolysis-a & $5.8694107$ & 1 & fitted \\
hydrolysis enthalpy increment & $0$ & J/mol & fitted \\
DEA / dispersion energy  /  Boltzmann constant & $269$ & K & fixed \\
DEA / segment count & $5.2048$ & 1 & fixed \\
DEA / segment diameter & $2.9477$ & \AA & fixed \\
DEA / solvation factor & $1$ & 1 & fixed \\
DEA / dielectric volume & $0.75748303$ & L/mol & fixed \\
DEA / dielectric temperature & $1155.7835$ & K & fixed \\
DEA / charge number & $0$ & $e$ & fixed \\
DEA / molar mass & $0.105137$ & kg/mol & fixed \\
DEA cation / Born diameter & $2.9477$ & \AA & fixed \\
DEA cation / Debye--H\"uckel diameter & $2.9477$ & \AA & fixed \\
DEA cation / dispersion energy  /  Boltzmann constant & $269$ & K & fixed \\
DEA cation / segment count & $5.2048$ & 1 & fixed \\
DEA cation / segment diameter & $2.9477$ & \AA & fixed \\
DEA cation / solvation factor & $1$ & 1 & fixed \\
DEA cation / charge number & $1$ & $e$ & fixed \\
DEA cation / molar mass & $0.106145$ & kg/mol & fixed \\
DEA cation / charged-chain indicator & $1$ & 1 & fixed \\
DEA carbamate / Born diameter & $6.1672093$ & \AA & fixed \\
DEA carbamate / Debye--H\"uckel diameter & $6.1672093$ & \AA & fixed \\
DEA carbamate / dispersion energy  /  Boltzmann constant & $183.00063$ & K & fixed \\
DEA carbamate / packing diameter & $5.4271442$ & \AA & fixed \\
DEA carbamate / segment count & $1$ & 1 & fixed \\
DEA carbamate / segment diameter & $6.1672093$ & \AA & fixed \\
DEA carbamate / solvation factor & $1$ & 1 & fixed \\
DEA carbamate / charge number & $-1$ & $e$ & fixed \\
DEA carbamate / molar mass & $0.148139$ & kg/mol & fixed \\
association / DEA / a / DEA / b / association energy  /  Boltzmann constant & $1517.6$ & K & fixed \\
association / DEA / a / DEA / b / volume & $0.1014$ & 1 & fixed \\
association / DEA cation / a / DEA cation / b / association energy  /  Boltzmann constant & $0$ & K & fixed \\
association / DEA cation / a / DEA cation / b / volume & $0.1014$ & 1 & fixed \\
association / DEA / a / multiplicity & $2$ & 1 & fixed \\
association / DEA / b / multiplicity & $2$ & 1 & fixed \\
association / DEA cation / a / multiplicity & $2$ & 1 & fixed \\
association / DEA cation / b / multiplicity & $2$ & 1 & fixed \\
DEA carbamate / hydronium / dispersion interaction & $0$ & 1 & fixed \\
DEA carbamate / carbon dioxide / dispersion interaction & $0$ & 1 & fixed \\
carbon dioxide / DEA / dispersion interaction & $0$ & 1 & fixed \\
carbon dioxide / DEA / reciprocal-temperature interaction slope & $0$ & K & fixed \\
carbon dioxide / DEA cation / dispersion interaction & $0$ & 1 & fixed \\
bicarbonate / DEA / dispersion interaction & $0$ & 1 & fixed \\
DEA carbamate / DEA / dispersion interaction & $0$ & 1 & fixed \\
carbonate / DEA / dispersion interaction & $0$ & 1 & fixed \\
DEA / hydronium / dispersion interaction & $0$ & 1 & fixed \\
DEA / hydroxide / dispersion interaction & $0$ & 1 & fixed \\
DEA / DEA cation / dispersion interaction & $0$ & 1 & fixed \\
DEA / water / dispersion interaction & $-0.063633523$ & 1 & fixed \\
DEA / water / reciprocal-temperature interaction slope & $0$ & K & fixed \\
bicarbonate / DEA cation / dispersion interaction & $0$ & 1 & fixed \\
DEA carbamate / DEA cation / dispersion interaction & $0$ & 1 & fixed \\
carbonate / DEA cation / dispersion interaction & $0$ & 1 & fixed \\
hydroxide / DEA cation / dispersion interaction & $0$ & 1 & fixed \\
DEA cation / water / reciprocal-temperature interaction slope & $0$ & K & fixed \\
DEA dissociation / a & $-9.2798784$ & 1 & fixed \\
DEA dissociation / b & $-4214.0761$ & K & fixed \\
DEA dissociation / c & $0$ & 1 & fixed \\
DEA dissociation / d & $0.0099612134$ & 1/K & fixed \\
DEA dissociation / e & $0$ & K$^2$ & fixed \\
DEA dissociation / f & $0$ & K$^{-2}$ & fixed \\
carbamate-hydrolysis / a & $5.8694107$ & 1 & fitted \\
carbamate-hydrolysis / b & $-2660.2065$ & K & fitted \\
carbamate-hydrolysis / c & $0$ & 1 & fixed \\
carbamate-hydrolysis / d & $0$ & 1/K & fixed \\
carbamate-hydrolysis / e & $0$ & K$^2$ & fixed \\
carbamate-hydrolysis / f & $0$ & K$^{-2}$ & fixed \\
\end{longtable}\endgroup


\section{Source conventions of the data}\label{sec:p2-s-observations}
% P:p2-s-observations
Tables~\ref{main-tab:p2-dea-data-current}, \ref{main-tab:p2-mdea-data-current} and \ref{main-tab:p2-amp-data-current} of the main text give the composition and loading ranges of each data set; the DEA and MDEA tables use the source's own basis, and the AMP table gives the unloaded AMP content in wt\% used in the calculation.
The 2.0 kmol/m$^3$ DEA concentration of Lal et al.\ is represented by unloaded mass fraction 0.2058 using solution density.
The species calculations for B\"ottinger et al.\ use 2.5 MPa because the source gives a 0.5--2.5 MPa cell-pressure interval without individual values.
All observations of Weiland et al.\ are kept, including the suspected source misprint identified in the archived observation file.
The \COtwo-free pressures of Kuranov et al.\ are read from a printed column that is 10 times the total pressure in MPa \cite[Table~2, p.~1960]{kuranovSolubilitySingleGases1996}.
The MDEA species calculations for B\"ottinger et al.\ are made at 0.5 and at 2.5 MPa, the ends of the reported cell-pressure interval, and both are reported.
The heats are read from Fig.~3 of Merkley et al.\ and the loading points from Fig.~4.
The assigned pressure is 1.121 MPa, or 1.466 MPa at 422.04 K.
The reading allowance for the loading point is 0.02 mol/mol.
Source measurement uncertainties are not agreement thresholds.

\section{Deviations inside and outside the measured ranges of the equilibrium constants}\label{sec:p2-s-ranges}
% P:p2-s-ranges
Tables~\ref{tab:p2-dea-ranges-current}, \ref{tab:p2-mdea-ranges-current} and \ref{tab:p2-amp-ranges-current} separate the deviations of each data set by its position relative to the measured temperature range of the carbamate hydrolysis constant (R4) and of the amine dissociation constant (R5); the main text gives the ranges and the rule that every observation is calculated.
MDEA has no R4, and the AMP R4 range is that of model-derived estimates.
Each range is assessed independently.
Counts include zero-valued species observations, while relative deviations use only nonzero observations.

\begingroup\small\setlength{\tabcolsep}{2pt}\begin{longtable}{@{}>{\raggedright\arraybackslash}p{0.33\linewidth}>{\raggedright\arraybackslash}p{0.12\linewidth}>{\raggedright\arraybackslash}p{0.14\linewidth}>{\raggedright\arraybackslash}p{0.08\linewidth}>{\raggedright\arraybackslash}p{0.08\linewidth}>{\raggedright\arraybackslash}p{0.14\linewidth}@{}}
\caption{DEA deviations separated by dissociation measurement interval (D) and hydrolysis source measurement interval (H). Binary comparisons do not use reactions and are omitted here.}\label{tab:p2-dea-ranges-current}\\
\toprule Source / quantity & Interval & Position & $n$ & $n_{rel}$ & AARD / \% \\ \midrule
\endfirsthead\toprule Source / quantity & Interval & Position & $n$ & $n_{rel}$ & AARD / \% \\ \midrule
\endhead\bottomrule\endfoot
B\"ottinger, carbamate & H & inside & 9 & 8 & $34.00$ \\
B\"ottinger, carbamate & H & outside & 16 & 15 & $72.66$ \\
B\"ottinger, carbamate & D & inside & 18 & 17 & $53.48$ \\
B\"ottinger, carbamate & D & outside & 7 & 6 & $75.43$ \\
B\"ottinger, DEA + cation & H & inside & 9 & 9 & $3.82$ \\
B\"ottinger, DEA + cation & H & outside & 16 & 16 & $9.52$ \\
B\"ottinger, DEA + cation & D & inside & 18 & 18 & $6.48$ \\
B\"ottinger, DEA + cation & D & outside & 7 & 7 & $10.01$ \\
B\"ottinger, bicarbonate & H & inside & 9 & 6 & $32.57$ \\
B\"ottinger, bicarbonate & H & outside & 16 & 13 & $33.31$ \\
B\"ottinger, bicarbonate & D & inside & 18 & 14 & $34.95$ \\
B\"ottinger, bicarbonate & D & outside & 7 & 5 & $27.84$ \\
Carson & H & outside & 50 & 50 & $0.95$ \\
Carson & D & inside & 50 & 50 & $0.95$ \\
Lal & H & inside & 17 & 17 & $56.30$ \\
Lal & H & outside & 27 & 27 & $42.68$ \\
Lal & D & inside & 17 & 17 & $56.30$ \\
Lal & D & outside & 27 & 27 & $42.68$ \\
Lawson & H & inside & 7 & 7 & $46.91$ \\
Lawson & H & outside & 23 & 23 & $22.85$ \\
Lawson & D & inside & 7 & 7 & $46.91$ \\
Lawson & D & outside & 23 & 23 & $22.85$ \\
Lawson, higher temperature & H & outside & 7 & 7 & $23.12$ \\
Lawson, higher temperature & D & outside & 7 & 7 & $23.12$ \\
Lawson, Table VIII & H & inside & 4 & 4 & $28.96$ \\
Lawson, Table VIII & H & outside & 6 & 6 & $43.67$ \\
Lawson, Table VIII & D & inside & 6 & 6 & $31.27$ \\
Lawson, Table VIII & D & outside & 4 & 4 & $47.57$ \\
Seo and Hong & H & inside & 5 & 5 & $31.13$ \\
Seo and Hong & H & outside & 11 & 11 & $109.41$ \\
Seo and Hong & D & inside & 5 & 5 & $31.13$ \\
Seo and Hong & D & outside & 11 & 11 & $109.41$ \\
Weiland & H & outside & 44 & 44 & $0.63$ \\
Weiland & D & inside & 44 & 44 & $0.63$ \\
\end{longtable}\endgroup

\begin{table}[htbp]\small\setlength{\tabcolsep}{2pt}\caption{Deviations of the present model from the MDEA data inside and outside the measured temperature ranges of the equilibrium constants.}\label{tab:p2-mdea-ranges-current}
\begin{tabular}{@{}>{\raggedright\arraybackslash}p{0.6\linewidth}>{\raggedright\arraybackslash}p{0.07\linewidth}>{\raggedright\arraybackslash}p{0.09\linewidth}>{\raggedright\arraybackslash}p{0.07\linewidth}>{\raggedright\arraybackslash}p{0.09\linewidth}@{}}
\toprule  & \multicolumn{2}{@{}l}{R5, inside} & \multicolumn{2}{@{}l}{R5, outside} \\ Data set & $N/N_{\mathrm{rel}}$ & AARD / \% & $N/N_{\mathrm{rel}}$ & AARD / \% \\ \midrule
Kuranov et al., total pressure & 37/37 & $11.49$ & 29/29 & $32.21$ \\
Jakobsen et al., species ratios & 16/16 & $13.62$ & --- & --- \\
Carson et al., integral heat & 40/40 & $5.54$ & --- & --- \\
Merkley et al., integral heat & 2/2 & $8.11$ & 2/2 & $4.82$ \\
Weiland et al., heat capacity & 33/33 & $2.29$ & --- & --- \\
Weiland et al., density & 33/33 & $0.22$ & --- & --- \\
Ermatchkov et al., $p_{\COtwo}$ & 69/69 & $15.52$ & 32/32 & $15.01$ \\
P\'erez-Salado Kamps et al., corrected temperatures, total pressure & 21/21 & $25.06$ & 7/7 & $25.99$ \\
B\"ottinger et al., MDEA + MDEAH$^+$, 0.5 MPa, mole fraction & 58/58 & $7.20$ & --- & --- \\
B\"ottinger et al., bicarbonate + carbonate, 0.5 MPa, mole fraction & 58/57 & $10.99$ & --- & --- \\
B\"ottinger et al., MDEA + MDEAH$^+$, 2.5 MPa, mole fraction & 58/58 & $7.20$ & --- & --- \\
B\"ottinger et al., bicarbonate + carbonate, 2.5 MPa, mole fraction & 58/57 & $10.99$ & --- & --- \\
Merkley et al., half the read loading, integral heat & 3/3 & $16.47$ & 3/3 & $10.09$ \\
Merkley et al., remaining states, integral heat & --- & --- & 1/1 & $0.65$ \\
Weiland et al., 60 wt\%, heat capacity & 6/6 & $9.24$ & --- & --- \\
Weiland et al., 60 wt\%, density & 11/11 & $1.92$ & --- & --- \\
Kuranov et al., 413 K, total pressure & --- & --- & 11/11 & $18.97$ \\
Zheng et al., integral heat & 6/6 & $17.21$ & --- & --- \\
Sedighi et al., total pressure & 15/15 & $71.37$ & --- & --- \\
Xiao et al., $p_{\COtwo}$ & 20/20 & $20.68$ & --- & --- \\
\bottomrule\end{tabular}
\par\smallskip{\footnotesize\raggedright $N/N_{\mathrm{rel}}$: calculated observations / observations with a defined relative deviation; ---: no observations. Sets that use no reaction are omitted.\par}\end{table}

\begin{table}[htbp]\small\setlength{\tabcolsep}{2pt}\caption{Deviations of the present model from the AMP data inside and outside the measured temperature ranges of the equilibrium constants.}\label{tab:p2-amp-ranges-current}
\begin{tabular}{@{}>{\raggedright\arraybackslash}p{0.28\linewidth}>{\raggedright\arraybackslash}p{0.07\linewidth}>{\raggedright\arraybackslash}p{0.09\linewidth}>{\raggedright\arraybackslash}p{0.07\linewidth}>{\raggedright\arraybackslash}p{0.09\linewidth}>{\raggedright\arraybackslash}p{0.07\linewidth}>{\raggedright\arraybackslash}p{0.09\linewidth}>{\raggedright\arraybackslash}p{0.07\linewidth}>{\raggedright\arraybackslash}p{0.09\linewidth}@{}}
\toprule  & \multicolumn{2}{@{}l}{R4, inside} & \multicolumn{2}{@{}l}{R4, outside} & \multicolumn{2}{@{}l}{R5, inside} & \multicolumn{2}{@{}l}{R5, outside} \\ Data set & $N/N_{\mathrm{rel}}$ & AARD / \% & $N/N_{\mathrm{rel}}$ & AARD / \% & $N/N_{\mathrm{rel}}$ & AARD / \% & $N/N_{\mathrm{rel}}$ & AARD / \% \\ \midrule
Hartono et al. (2021), 313--353 K, total pressure & 32/32 & $11.14$ & 15/15 & $8.71$ & 47/47 & $10.36$ & --- & --- \\
Ciftja et al., carbamate fraction & 16/16 & $23.57$ & 41/41 & $28.95$ & 57/57 & $27.44$ & --- & --- \\
Hartono et al. (2021), 373--393 K, total pressure & 11/11 & $26.42$ & 7/7 & $30.82$ & --- & --- & 18/18 & $28.13$ \\
Kundu et al., $p_{\COtwo}$ & 11/11 & $17.88$ & 40/40 & $51.27$ & 51/51 & $44.07$ & --- & --- \\
Tong et al., total pressure & 17/17 & $29.17$ & 23/23 & $27.42$ & 17/17 & $29.17$ & 23/23 & $27.42$ \\
\bottomrule\end{tabular}
\par\smallskip{\footnotesize\raggedright $N/N_{\mathrm{rel}}$: calculated observations / observations with a defined relative deviation; ---: no observations. Sets that use no reaction are omitted.\par}\end{table}


% P:p2-s-ranges-alternative
For the stability-constrained AMP alternative, the pressure AARDs inside and outside R4, respectively, are 19.78 and 20.60\,\% for the fitted Hartono data, 21.55 and 28.92\,\% for the higher-temperature Hartono data, 30.91 and 111.59\,\% for Kundu, and 52.99 and 24.73\,\% for Tong.
All fitted Hartono and Kundu pressures lie inside R5 (20.04 and 94.19\,\%), and all higher-temperature Hartono pressures lie outside it (24.42\,\%).
For Tong, the R5 groups coincide with the R4 groups.

\section{Published comparisons on the same data}\label{sec:p2-s-published}
% P:p2-s-published
Lawson and Garst report scatter against smoothed curves, without a same-observation scalar thermodynamic-model deviation \cite[p.~26, Results and Accuracy of data]{lawsonGasSweeteningData1976}.
Seo and Hong's graphical smoothed comparison is for AMP, not the DEA subset used here \cite[Table~1, Results and Discussion, p.~259]{seoSolubilitiesCarbonDioxide1996}.
Lal et al. compare the Kent--Eisenberg treatment graphically with observations not used to fit those earlier constants; no scalar deviation on the present subset is reported \cite[pp.~683--684, Eqs.~9--13, Fig.~3]{lalSolubilityH2SCO21985}.
The infinite-dilution comparison of Carson et al.\ concerns a different quantity from the finite-loading integral heats fitted here \cite[Tables~2, 5--6, pp.~1289--1295]{carsonEnthalpySolutionCarbon2000}.
Horstmann et al.\ fit binary vapor--liquid equilibrium and excess enthalpy simultaneously; the same heat observations are comparison only here \cite[pp.~1496--1500, Tables~6--7, Fig.~4(b)]{horstmannPhaseEquilibriumExcess2002}.
The source's excess-enthalpy uncertainty expression has no absolute-value bars on its enthalpy-dependent term despite negative measured heats; it is not used to compute a negative uncertainty.

\section{Stability penalty and retreat procedure}\label{sec:p2-s-retreat}
% P:p2-s-penalty
The main text defines the phase-stability search and states when the stability penalty and the retreat apply; this section gives both procedures in full.
Where a liquid descent is found, the fit is repeated with Part I's stability penalty, which adds half the squared, scaled most negative tangent-plane distance of each penalized state to the cost \cite[Sect.~3.4]{PolleyPartI}; the penalized states are compressed-liquid copies of the whole searched liquid-state set, including measured-state liquids and grid states, at both pressure factors, with identical copies counted once.
Later fitting stages keep the compositions of liquid descents already found as supplied trials, so a detected descent is tested again.
If neither the penalty nor the retreat below removes the split, the converged unpenalized solution is reported, with its detected split counts as lower bounds.

% P:p2-s-retreat
If the penalized solution has a detected descent or an unavailable search, a retreat is attempted from a reference parameter set with no detected descent and no unavailable search.
A point is admissible for retreat only when its fitted predictions can be recalculated, all required searches are available, and neither an optimized trial nor a supplied composition indicates a descent.
Points are interpolated linearly between the reference set at $t=0$ and the penalized solution at $t=1$; bisection narrows the bracket to at most $0.01/2^{10}$, the resolution of Part I's 0.01 steps followed by 10 bisections \cite[Sect.~3.4]{PolleyPartI}.
Every point is searched with a copy of the trial set fixed after the 2 endpoint searches, plus only the descent compositions found at that point; its search is repeated until no new descent composition is found.
Thus trials discovered at other interpolated points do not change the decision at a given $t$.
Bisection assumes a single transition from admissible to inadmissible as $t$ increases.
The selected point is checked again; when it lies below $t=1$, the upper end of the final bracket is also checked, and all observed reversals of admissibility are reported.
No candidate is returned if the reference point fails these checks.
An interpolated point is not a converged fit.
If the selected point is the fitted endpoint at $t=1$, its convergence result applies.

% P:p2-s-penalty-amines
The DEA fit used neither the penalty nor the retreat.
For MDEA, the penalized states are 78 \COtwo-free MDEA--water liquids at 1.466 MPa: the temperatures 393, 403, 413, 413.15, 422.04 and 423 K, each at MDEA mole fractions of 0.01, 0.02, 0.04, 0.06, 0.07, 0.08, 0.10, 0.12, 0.14, 0.16, 0.18, 0.20 and 0.24.
The penalty scales are 0.01 and then 0.003, and the retreat interpolates between each initial parameter set, which shows no descent at these states, and its penalized solution.
The accepted points lie at $t=0.8690$ and $t=0.8688$ for the 2 initial sets; at each, all 78 searches complete without a descent, and the next point above has 1 unavailable search.
No penalized fit was made with the compressed-liquid copies of the loaded MDEA states; one fitting stage on the full set of 11542 penalized states was projected from measured evaluation times to take 14.1--20.2 h.
For AMP, the penalized states are the 862 compressed-liquid copies of the searched states, the penalty scales are 0.01 and 0.003, and the accepted retreat point lies at $t=0.5699$; the next point above has 1 unavailable search.

\section{Phase-search conditions and reproducibility}\label{sec:p2-s-numerical}
% P:p2-s-numerical
The electronic data contain the parameter and input files with their checksums, the calculated value or reason for unavailability for each observation, the conditional information matrices and the definition of every searched state.
Trial phases exclude species that cannot form from the feed at the searched state.
At the second observation in Table VIII of Lawson and Garst (298.15 K, 54.5 wt\% DEA, loading 0.602 mol/mol), the first trial converged but its reconstructed composition did not preserve the required species.
The phase search was unavailable; the pressure observation was calculated.
Recalculating the 80 fitted predictions from the archived DEA parameters reproduces the selected solution, which checks the numerical calculation rather than the physical description.

% P:p2-s-numerical-mdea
Of the 1039 MDEA search requests with incomplete results, 810 are grid requests without a calculated bubble-pressure reference.
The MDEA liquid--liquid split at 363.15 K has a tangent-plane distance of $-0.00262$ in units of $RT$.
The 27 MDEA calculations with a descent toward a vapor are distinct calculations at conditions assigned to measurements, whose assigned pressure is 0.1733--0.9990 times the calculated bubble pressure: 4 densities of Sobrino et al., 11 species states of B\"ottinger et al.\ at 0.5 MPa and 2 at 2.5 MPa, 5 heats of Zheng et al.\ and 5 heats of Merkley et al.
At 3 of the 11 species states at 0.5 MPa the descent is shown only by a supplied trial composition.

% P:p2-s-numerical-amp
For AMP, each of the 90 descents was classified by solving for the density of the trial phase at the pressure of the parent liquid; every parent and trial phase is liquid.
A liquid--liquid split is counted as confirmed when a separate equilibrium calculation at fixed feed, temperature and pressure returns 2 distinct liquids with positive amounts and a Gibbs-energy decrease larger than its estimated numerical error.
The largest decrease occurs at 393.15 K, 30.0 wt\% AMP and a loading of 0.700, at 1.215 MPa, where the 2 liquid densities are 44685 and 31704 mol/m$^3$.

% P:p2-s-software
The calculations for the 3 amines used one build of the \texttt{epcsaft} package, version 1.3.1, identified by the software index in the archive deposit \textit{[Zenodo DOI supplied in the final submission]}.
